\documentclass{article}

% ── Encoding & fonts ─────────────────────────────────────────────────────────
\usepackage[utf8]{inputenc}
\usepackage[T1]{fontenc}
\usepackage{lmodern}

% ── Page geometry ─────────────────────────────────────────────────────────────
\usepackage[a4paper, margin=2.5cm]{geometry}

% ── Hyperlinks ────────────────────────────────────────────────────────────────
\usepackage[hidelinks, unicode]{hyperref}

% ── Source code listings ─────────────────────────────────────────────────────
\usepackage{listings}
\lstset{
  basicstyle=\ttfamily\small,
  breaklines=true,
  keepspaces=true,
  columns=flexible,
  frame=single,
  numbers=left,
  numberstyle=\tiny,
}

% ── List customisation ───────────────────────────────────────────────────────
\usepackage{enumitem}

% ── Strikethrough ────────────────────────────────────────────────────────────
\usepackage[normalem]{ulem}

% ── Graphics ─────────────────────────────────────────────────────────────────
\usepackage{graphicx}

% ── Tables ───────────────────────────────────────────────────────────────────
\usepackage{booktabs}
\usepackage{tabularx}

\begin{document}
\section{MarkTeX Technical Reference}

\subsection{Overview}

MarkTeX converts Markdown to LaTeX via a typed AST. All source files are UTF-8; output targets \texttt{pdflatex}, \texttt{xelatex}, or \texttt{lualatex} with the packages listed in Table~\ref{tab:tab_extensions}.

\subsubsection{Design Goals}

\begin{enumerate}
  \item \textbf{Correctness}: special characters are always escaped (see Table~\ref{tab:tab_escaping}).
  \item \textbf{Extensibility}: new constructs follow the five-step pattern described in~\cite{knuth1986}.
  \item \textbf{Transparency}: the AST is the only boundary between parsing and generation, as specified in~\cite{latex2023guide}.
\end{enumerate}


\noindent\rule{\linewidth}{0.4pt}

\subsection{Definition Blocks}

A \texttt{---…---} block defines all reference keys and document metadata in one place. Multiple entry types may coexist in the same block:

\begin{lstlisting}[language=markdown]
---
C#01:knuth1986
F#01:my_figure
T#01:my_table
Author: Ada Lovelace
Title: My Document
MakeTOC
---
\end{lstlisting}

The block is pre-scanned before parsing begins, so keys defined at the bottom of a file are available at the top. A \texttt{---} line that is not followed by a recognised entry type is treated as a thematic break (\texttt{\textbackslash{}noindent\textbackslash{}rule\{\textbackslash{}linewidth\}\{0.4pt\}}).


\noindent\rule{\linewidth}{0.4pt}

\subsection{AST Node Types}

Table~\ref{tab:my_table} lists all block and inline node types with their LaTeX counterparts.

\begin{table}[H]
\centering
\caption{Complete node-type reference. Leaf nodes have no children. See Figure~\ref{fig:my_figure} for the hierarchy.}
\begin{tabular}{lclc}
\toprule
Node type & Category & LaTeX output & Leaf \\
\midrule
Document & block & (root) & no \\
Heading & block & \texttt{\textbackslash{}section}, \texttt{\textbackslash{}subsection}, … & no \\
Paragraph & block & text + \texttt{\textbackslash{}n\textbackslash{}n} & no \\
BulletList & block & \texttt{\textbackslash{}begin\{itemize\}} & no \\
OrderedList & block & \texttt{\textbackslash{}begin\{enumerate\}} & no \\
ListItem & block & \texttt{\textbackslash{}item} & no \\
BlockQuote & block & \texttt{\textbackslash{}begin\{quote\}} & no \\
FencedCode & block & \texttt{\textbackslash{}begin\{lstlisting\}} & yes \\
IndentedCode & block & \texttt{\textbackslash{}begin\{verbatim\}} & yes \\
ThematicBreak & block & \texttt{\textbackslash{}noindent\textbackslash{}rule\{…\}} & yes \\
Table & block & \texttt{\textbackslash{}begin\{tabular\}} & no \\
TableBlock & block & \texttt{\textbackslash{}begin\{table\}…\textbackslash{}end\{table\}} & no \\
DefinitionBlock & block & (invisible) & yes \\
FigureBlock & block & (invisible) & yes \\
Text & inline & escaped text & yes \\
Strong & inline & \texttt{\textbackslash{}textbf\{\}} & no \\
Emphasis & inline & \texttt{\textbackslash{}textit\{\}} & no \\
StrongEmphasis & inline & \texttt{\textbackslash{}textbf\{\textbackslash{}textit\{\}\}} & no \\
CodeSpan & inline & \texttt{\textbackslash{}texttt\{\}} & yes \\
Strikethrough & inline & \texttt{\textbackslash{}sout\{\}} & no \\
Link & inline & \texttt{\textbackslash{}href\{\}\{\}} & no \\
Image & inline & \texttt{\textbackslash{}includegraphics\{\}} & no \\
FigureImage & inline & \texttt{\textbackslash{}begin\{figure\}…\textbackslash{}end\{figure\}} & yes \\
CitationRef & inline & \texttt{\textasciitilde{}\textbackslash{}cite\{\}} & yes \\
FigureRef & inline & \texttt{\textasciitilde{}\textbackslash{}ref\{fig:\}} & yes \\
TableRef & inline & \texttt{\textasciitilde{}\textbackslash{}ref\{tab:\}} & yes \\
\bottomrule
\end{tabular}
\label{tab:my\_table}
\end{table}


\noindent\rule{\linewidth}{0.4pt}

\subsection{Special Character Escaping}

Table~\ref{tab:tab_escaping} shows the ten LaTeX special characters and their escaped forms.

\begin{table}[ht]
\centering
\caption{LaTeX special-character escape table. Escaping is \textbf{not} applied inside \texttt{lstlisting} or \texttt{verbatim} environments, nor inside URL arguments to \texttt{\textbackslash{}href\{\}}. See~\cite{latex2023guide} for details.}
\begin{tabular}{lll}
\toprule
Character & Name & LaTeX output \\
\midrule
\texttt{\textbackslash{}} & backslash & \texttt{\textbackslash{}textbackslash\{\}} \\
\texttt{\{} & left brace & \texttt{\textbackslash{}\{} \\
\texttt{\}} & right brace & \texttt{\textbackslash{}\}} \\
\texttt{\$} & dollar & \texttt{\textbackslash{}\$} \\
\texttt{\&} & ampersand & \texttt{\textbackslash{}\&} \\
\texttt{\#} & hash & \texttt{\textbackslash{}\#} \\
\texttt{\%} & percent & \texttt{\textbackslash{}\%} \\
\texttt{\_} & underscore & \texttt{\textbackslash{}\_} \\
\texttt{\textasciicircum{}} & caret & \texttt{\textbackslash{}textasciicircum\{\}} \\
\texttt{\textasciitilde{}} & tilde & \texttt{\textbackslash{}textasciitilde\{\}} \\
\bottomrule
\end{tabular}
\label{tab:tab\_escaping}
\end{table}


\noindent\rule{\linewidth}{0.4pt}

\subsection{Figures}

Figure~\ref{fig:my_figure} shows the overall architecture. Figure~\ref{fig:fig_flow} shows the data flow in detail, and Figure~\ref{fig:fig_output} shows a sample output page.

\begin{figure}[ht]
  \centering
  \includegraphics[width=1.0\linewidth]{figures/arch.png}
  \caption{System architecture overview. See~\cite{knuth1986} for the Go implementation.}
  \label{fig:my\_figure}
\end{figure}

\begin{figure}[h]
  \centering
  \includegraphics[width=0.8\linewidth]{figures/flow.png}
  \caption{Data-flow diagram. The pre-scan phase resolves~\cite{knuth1986},~\ref{fig:my_figure}, and~\ref{tab:my_table} references before parsing.}
  \label{fig:fig\_flow}
\end{figure}

\begin{figure}[hb]
  \centering
  \includegraphics[width=0.6\linewidth]{figures/output.png}
  \caption{Sample LaTeX output. The document uses packages listed in Table~\ref{tab:tab_extensions}.}
  \label{fig:fig\_output}
\end{figure}


\noindent\rule{\linewidth}{0.4pt}

\subsection{Extensions Reference}

Table~\ref{tab:tab_extensions} lists all available extensions, the LaTeX packages they require, and the \texttt{-ext} token used on the command line.

\begin{table}[ht]
\centering
\caption{Extension reference. Packages marked * are only included when actually used. See Figure~\ref{fig:my_figure} for where each extension fits in the pipeline.}
\begin{tabular}{lll}
\toprule
\texttt{-ext} token & Feature & Required package \\
\midrule
\texttt{tables} & GFM pipe tables & \texttt{booktabs} * \\
\texttt{strikethrough} & \texttt{\textasciitilde{}\textasciitilde{}text\textasciitilde{}\textasciitilde{}} & \texttt{ulem} * \\
\texttt{citations} & \texttt{[C\#key]} → \texttt{\textbackslash{}cite\{\}} & — \\
\texttt{figures} & \texttt{[F\#key]} → \texttt{\textbackslash{}ref\{fig:\}} & \texttt{graphicx} * \\
\texttt{tableFloat} & \texttt{[T\#key]} → \texttt{\textbackslash{}ref\{tab:\}} & \texttt{booktabs} * \\
\texttt{documentMeta} & \texttt{Author:}, \texttt{MakeTOC}, … & — \\
\texttt{all} & All of the above & as needed \\
\bottomrule
\end{tabular}
\label{tab:tab\_extensions}
\end{table}


\noindent\rule{\linewidth}{0.4pt}

\subsection{Inline Formatting Reference}

The following examples show all inline constructs in context. A \sout{deprecated} feature is shown for completeness. The \texttt{\textbackslash{}texttt\{monospace\}} style is used for \texttt{inline code}. \textbf{Bold text} and \textit{italic text} and \textbf{\textit{bold-italic text}} are all supported. Links render as \texttt{\textbackslash{}href\{url\}\{text\}} — for example \href{https://go.dev/ref/spec}{the Go specification}.

Hard line break example (two trailing spaces): First line. Second line after hard break.

\begin{quote}
Blockquotes nest the content in \texttt{\textbackslash{}begin\{quote\}…\textbackslash{}end\{quote\}}. A citation~\cite{latex2023guide} and a figure reference~\ref{fig:my_figure} both resolve correctly inside a blockquote.

\end{quote}


\noindent\rule{\linewidth}{0.4pt}

\subsection{Standalone Output}

When \texttt{-standalone} is active, the preamble includes only the packages the document actually needs. For a document with code listings, tables, images, and strikethrough, the generated preamble is:

\begin{lstlisting}[language=latex]
\documentclass{article}
\usepackage[utf8]{inputenc}
\usepackage[T1]{fontenc}
\usepackage{lmodern}
\usepackage[a4paper, margin=2.5cm]{geometry}
\usepackage[hidelinks, unicode]{hyperref}
\usepackage{listings}
\usepackage{enumitem}
\usepackage[normalem]{ulem}
\usepackage{graphicx}
\usepackage{booktabs}
\usepackage{tabularx}
\end{lstlisting}

Document metadata (\texttt{Author:}, \texttt{Title:}, etc.) is placed immediately before \texttt{\textbackslash{}begin\{document\}}, and front-matter commands (\texttt{\textbackslash{}maketitle}, \texttt{\textbackslash{}tableofcontents}, etc.) are placed immediately after it.


\end{document}
